\section{ファイル構成}

\subsection{ディレクトリツリー}

\begin{lstlisting}[basicstyle=\ttfamily\footnotesize,language={},numbers=none]
font_to_bin/
├── public/
│   ├── fonts/
│   │   ├── DotGothic16-Regular.ttf
│   │   └── misaki_gothic_2nd.ttf
│   ├── favicon.svg
│   └── manifest.webmanifest
├── src/
│   ├── core/
│   │   ├── types.ts       … 型定義（AppState, Glyph, Format 他）
│   │   ├── fonts.ts       … 同梱フォント定義 + カスタム登録API
│   │   ├── rasterize.ts   … Canvas によるグリフ2値化
│   │   ├── encode.ts      … ビットパック・数値フォーマット
│   │   ├── format.ts      … 言語別出力文字列生成
│   │   ├── share.ts       … URLクエリパラメータ共有
│   │   ├── presets.ts     … 用途別プリセット(FormatOptions雛形)
│   │   └── defaults.ts    … 既定値
│   ├── components/
│   │   ├── ui/            … shadcn/ui プリミティブ群
│   │   ├── Header.tsx
│   │   ├── EasyPanel.tsx        … かんたんモード(3ステップ)
│   │   ├── InputPanel.tsx
│   │   ├── FormatPanel.tsx
│   │   ├── PreviewPanel.tsx     … 横スクロール帯プレビュー
│   │   ├── GlyphPreview.tsx     … 補助線対応のグリフ描画
│   │   ├── OutputPanel.tsx
│   │   ├── PixelEditor.tsx
│   │   └── FreeCanvasPanel.tsx  … 自由モードのキャンバス
│   ├── hooks/
│   │   ├── useAppState.ts … useReducer ラッパ
│   │   ├── useGlyphs.ts   … メモ化ラスタライザ
│   │   ├── usePersist.ts  … LocalStorage自動保存
│   │   ├── useMode.ts     … かんたん/詳細モード切替
│   │   └── useTheme.ts    … ダーク/ライト
│   ├── lib/utils.ts       … cn() ヘルパ
│   ├── App.tsx            … 3カラムレイアウト + モーダル制御
│   ├── main.tsx           … エントリポイント
│   └── index.css          … Tailwind ベースCSS
├── .github/workflows/deploy.yml   … GitHub Pages 自動デプロイ
├── doc/                           … LaTeX 設計仕様書
├── legacy/                        … 旧Python版（参考保存）
├── index.html
├── package.json
├── tsconfig.json
├── tsconfig.node.json
├── vite.config.ts
├── tailwind.config.js
├── postcss.config.js
└── README.md
\end{lstlisting}

\subsection{core 層}

\begin{itemize}[leftmargin=*,itemsep=2pt]
  \item \texttt{src/core/types.ts}\,---\,\texttt{AppState}, \texttt{Glyph},
        \texttt{FormatOptions} など全型定義。UI と core の契約。
  \item \texttt{src/core/fonts.ts}\,---\,\texttt{BUILTIN\_FONTS} 配列、
        \texttt{loadFont}, \texttt{registerCustomFont}。
        FontFace API でカスタム TTF/OTF を登録。
  \item \texttt{src/core/rasterize.ts}\,---\,\texttt{rasterizeChar(char, opts)}
        $\rightarrow$ \texttt{Glyph}。Canvas へ1文字描画し
        \texttt{getImageData} でアルファ値を閾値2値化。
  \item \texttt{src/core/encode.ts}\,---\,\texttt{packRow} / \texttt{packColumn}
        （MSB/LSB ビットパック）、\texttt{formatNumber}（bin/hex/dec）。
  \item \texttt{src/core/format.ts}\,---\,\texttt{formatOutput}。
        出力言語に応じて \texttt{formatC} / \texttt{formatPython}
        / \texttt{formatJS} / \texttt{formatRust} / \texttt{formatGo}
        / \texttt{formatJSON} / \texttt{formatMarkdown}
        / \texttt{formatSymbols} / \texttt{formatPlain} を呼び分け。
  \item \texttt{src/core/share.ts}\,---\,\texttt{encodeStateToURL}
        / \texttt{decodeStateFromURL}。
        \texttt{AppState} を base64url 化して \texttt{?s=...} に載せる。
  \item \texttt{src/core/presets.ts}\,---\,\texttt{PRESETS} 配列
        （Arduino/C言語, C言語2D, Python リスト, MicroPython bytes,
        JSON, 記号アートの6種）と \texttt{detectPreset()}
        （現在の \texttt{FormatOptions} から該当プリセットIDを推定）。
        かんたんモードでワンタッチ切替に使用。
  \item \texttt{src/core/defaults.ts}\,---\,\texttt{DEFAULT\_STATE}
        （text=\texttt{"ABC"}, fontId=\texttt{dotgothic16},
        16$\times$16, threshold=128）。
\end{itemize}

\subsection{hooks 層}

\begin{itemize}[leftmargin=*,itemsep=2pt]
  \item \texttt{src/hooks/useAppState.ts}\,---\,\texttt{useReducer}。
        Action は \texttt{setText}, \texttt{setSize},
        \texttt{setFormat}, \texttt{setOverride} のほか、
        プレビュー補助線用 \texttt{setPreviewGridStep}、
        自由モード用 \texttt{setFreeSize} / \texttt{setFreeMatrix} /
        \texttt{setFreePixel} / \texttt{setFreeGridStep} /
        \texttt{clearFree} を提供する。
  \item \texttt{src/hooks/useGlyphs.ts}\,---\,text / font / size
        変更時のみ再ラスタライズするメモ化。overrides をマージ。
  \item \texttt{src/hooks/usePersist.ts}\,---\,LocalStorage キー
        \texttt{font-to-bin.state.v1}。
  \item \texttt{src/hooks/useMode.ts}\,---\,かんたん/詳細/自由モード切替。
        LocalStorage キー \texttt{font\_to\_bin\_mode\_v1}。値は
        \texttt{easy} / \texttt{advanced} / \texttt{free} のいずれか。
  \item \texttt{src/hooks/useTheme.ts}\,---\,light / dark 切替、
        \texttt{font-to-bin.theme} キー。
\end{itemize}

\subsection{UI 層}

\begin{itemize}[leftmargin=*,itemsep=2pt]
  \item \texttt{src/components/Header.tsx}\,---\,タイトル、
        \textbf{かんたん/詳細モード切替}のセグメントタブ、
        テーマ切替、URL 共有、GitHub リンク。
  \item \texttt{src/components/EasyPanel.tsx}\,---\,
        かんたんモード用の3ステップUI。
        \textbf{ステップ1}: 変換したい文字。
        \textbf{ステップ2}: フォント + ドットサイズクイックボタン
        (8/12/16/24/32px) + 太字トグル。
        \textbf{ステップ3}: 用途プリセットを6種から1クリック選択
        （\texttt{presets.ts} の \texttt{PRESETS} を使用）。
  \item \texttt{src/components/InputPanel.tsx}\,---\,詳細モードの左上。
        テキスト、フォント選択、幅/高さ、描画オフセット、太字、閾値。
  \item \texttt{src/components/FormatPanel.tsx}\,---\,出力言語、
        構造、基数、MSB/LSB、0/1 反転、変数名、データ型。
  \item \texttt{src/components/PreviewPanel.tsx}\,---\,グリフ一覧を
        \textbf{横スクロール帯}で表示。補助線間隔セレクタ、
        ズーム（50--400\%）を備える。
        文字クリックでピクセルエディタを開く。
  \item \texttt{src/components/GlyphPreview.tsx}\,---\,単一グリフの
        Canvas タイル描画。\texttt{gridStep} プロパティで N ドット毎の
        強調補助線を描画。
  \item \texttt{src/components/OutputPanel.tsx}\,---\,生成結果の
        表示 + コピー + ダウンロード。
  \item \texttt{src/components/PixelEditor.tsx}\,---\,モーダル。
        ペン/消しゴム/トグル/反転/回転/クリア/ズームに加え、
        補助線間隔の切替に対応。
  \item \texttt{src/components/FreeCanvasPanel.tsx}\,---\,自由モード専用。
        W$\times$H(1--256)指定、スクロール可能なキャンバス、補助線、
        ズーム、ペン／消しゴム／トグル、反転・回転・全消去。
        編集結果は \texttt{AppState.free.matrix} に
        \texttt{setFreePixel}（1セル単位）で反映され、
        出力側では単一グリフとして扱う。
  \item \texttt{src/components/ui/}\,---\,Radix UI ベースの
        shadcn/ui 汎用プリミティブ群。
\end{itemize}

\subsection{その他}

\begin{itemize}[leftmargin=*,itemsep=2pt]
  \item \texttt{.github/workflows/deploy.yml}\,---\,Node 20 で
        \texttt{npm ci} $\rightarrow$ \texttt{build} $\rightarrow$
        Pages にアップロード。
  \item \texttt{public/fonts/*.ttf}\,---\,同梱フォント
        （DotGothic16, Misaki Gothic 2nd）。
  \item \texttt{legacy/}\,---\,旧 Python/Tkinter 実装
        （参照用、ビルドからは除外）。
\end{itemize}

\newpage
